% =============================================================================
% DEFINITIONS (2-3 pages)
%
% Spec requirements:
%   Establish minimal formal vocabulary for the propositions.
%   Definitions 0-6 plus implementation notes.
%
% Key design decision (resolved in spec):
%   Candidate B (capabilities) as formal target.
%   Candidate C (experiences) as motivating justification.
%   "GFM does not claim to know what experiences agents want. It maximizes
%    the capabilities that let agents choose for themselves."
%
% Reference material:
%   - Predraft Part III opening and Equation 5
%   - Blog: free.md (informal goal frontier definitions)
%   - Klyubin et al. (2005), Sen (1999)
% =============================================================================

\section{Definitions}
\label{sec:definitions}

% ---------------------------------------------------------------------------
% Definition 0: What Is a Goal?
% The paper's core philosophical move: capabilities (B) as formal proxy,
% experiences (C) as justification. Make the B-as-proxy-for-C relationship
% explicit here, and return to it in Section 8 (Doll Problem, wireheading).
% ---------------------------------------------------------------------------

\begin{definition}[Goal]
\label{def:goal}
% TODO: Formal definition adopting Candidate B (capabilities)
% with Candidate C (experiences) as motivating justification
\end{definition}

% ---------------------------------------------------------------------------
% Definition 1: Capability Space, Agent, Individual Capability Set
% G is the universe of all possible capabilities.
% Agent a_k has individual capability set G_k^ind \subseteq G.
% ---------------------------------------------------------------------------

\begin{definition}[Capability Space, Agent, and Individual Capability Set]
\label{def:capability_space}
The \emph{capability space} $\G$ is the universe of all possible capabilities---everything any agent or coalition of agents could do or be. An \emph{agent} $a_k$ is an entity with an \emph{individual capability set} $\Gind{k} \subseteq \G$: the capabilities $a_k$ can realize unilaterally given its current resources, embodiment, and environment.

A capability $g \in \Gind{k}$ represents something $a_k$ \emph{can} do or be, not necessarily something it \emph{is} doing or \emph{will} do.
\end{definition}

% ---------------------------------------------------------------------------
% Definition 2: Joint Goal Space
% G = union of G_k^ind plus G^coop
% Canonical measure: vol(G)
% ---------------------------------------------------------------------------

\begin{definition}[Joint Goal Space]
\label{def:joint_goal_space}
The \emph{joint goal space} $G$ for a population of agents $\{a_1, \ldots, a_n\}$ is the set of all capabilities the population can realize through any combination of individual and coordinated action:
\begin{equation}
\label{eq:joint_goal_space}
G = \bigcup_k \Gind{k} \;\cup\; \Gcoop
\end{equation}
where $\Gcoop \subseteq \G$ contains capabilities achievable only through multi-agent cooperation. By construction, $\Gind{k} \subseteq G \subseteq \G$ for all $k$.

The \emph{frontier} $\partial G$ is the boundary of $G$. The \emph{canonical measure} is $\vol(G)$: the volume of the jointly achievable capability space. This is the quantity we optimize.
\end{definition}

% ---------------------------------------------------------------------------
% Definition 3: Goal-Frontier Maximizer (GFM)
% ---------------------------------------------------------------------------

\begin{definition}[Goal-Frontier Maximizer]
\label{def:gfm}
A \emph{goal-frontier maximizer} (GFM) actor selects actions $\pi$ to maximize the expected change in the volume of the joint goal space:
\begin{equation}
\label{eq:gfm_objective}
\pi^* = \arg\max_\pi \;\mathbb{E}\bigl[\Delta \vol(G) \mid \pi\bigr]
\end{equation}
\end{definition}

% ---------------------------------------------------------------------------
% Definition 4: Contraction and Expansion
% ---------------------------------------------------------------------------

\begin{definition}[Contraction and Expansion]
\label{def:contraction_expansion}
An action $\pi$ is \emph{expanding} if $\mathbb{E}[\Delta \vol(G) \mid \pi] > 0$ and \emph{contracting} if $\mathbb{E}[\Delta \vol(G) \mid \pi] < 0$. Contraction includes both direct contraction (destroying an agent's capabilities) and indirect contraction (restricting agency through coercion, deception, or rigid rules).
\end{definition}

% ---------------------------------------------------------------------------
% Definition 5: Social Objective
% Critical assumptions: additivity under independence, tradeoff acknowledgment
% ---------------------------------------------------------------------------

\begin{definition}[Social Objective]
\label{def:social_objective}
The social objective $V(G) = \vol(G)$ is a function of the joint goal space, not a sum of individual utilities.

\emph{Additivity under independence:} When agents' individual capability sets are independent (non-overlapping, non-interacting), $V(G) \geq \sum_k \vol(\Gind{k})$. Cooperative capabilities $\Gcoop$ contribute additional volume.

\emph{Tradeoff acknowledgment:} When expanding $G_j$ requires contracting $G_k$, the GFM actor must evaluate the \emph{net} change in $\vol(G)$. The framework makes the tradeoff explicit and measurable; it does not claim all tradeoffs are resolvable.
\end{definition}

% ---------------------------------------------------------------------------
% Definition 6: Observable Goal Model
% ---------------------------------------------------------------------------

\begin{definition}[Observable Goal Model]
\label{def:observable_goal_model}
A GFM actor maintains a model $M_k(G)$ of each observed agent's goals, inferred from communication, behavior, and environmental observation. Each model has an associated trust factor $T_k \in [0, 1]$ and confidence $C_k$.
\end{definition}
